% Physical unit 032; accepted Session 03 after fresh 400-dpi authority replay.
\GordanUnitStart{032}{printed page 27}
\GordanTopFolio{27}
\begingroup\small
from the formulae:
\[
\begin{array}{rcll}
 \sqrt{1-e^2}\,\tg B_0&=&\tg l_0, &\text{Cap. I \S\ 7}\vspace{0.2em}\\
 {2\over\pi}{\Sigma(l_0)\over b}&=&1+o\sin^2l_0-{3\over4}o^2\sin^4l_0+{5\over4}o^3\sin^6l_0\ldots, &\text{Cap. III \S\ 2}\vspace{0.2em}\\
 {a\over b}\left\{1-{2\over\pi}\mathcal A(l_0)\right\}&=&2p\cos l_0+\left\{1-2p+{25\over8}p^2\right\}p^2\cos^3l_0\ldots, &\text{Cap. III \S\ 4}\vspace{0.2em}\\
 x&=&e\sin B_0, &\text{Cap. II \S\ 9}\vspace{0.2em}\\
 x'&=&{\cos B_0\over \cos l_0}. &
\end{array}
\]

\sect{4}

It remains to compute the constants
\[
 q\qquad x\qquad E\qquad v.
\]
For this purpose it is useful to apply the following formula:
\[
 \sqrt{x'}={1-2q+2q^4\cdots\over 1+2q+2q^4\cdots}\qquad \text{fund. pag. 184.}
\]
Since $q$ is a very small quantity, which indeed, as regards the earth, never exceeds the number \gq{0,0004}, $q^4$ may be wholly neglected.

Therefore there is:
\[
 \sqrt{x'}={1-2q\over 1+2q}
\]
or:
\[
 2q={1-\sqrt{x'}\over 1+\sqrt{x'}}
 ={\sqrt{\cos l_0}-\sqrt{\cos B_0}\over \sqrt{\cos l_0}+\sqrt{\cos B_0}},
\]
from which $q$ can easily be computed.

In this matter, since $x'$ is sometimes not accurately defined, we grant that it may happen that $q$ cannot be computed to the extreme digits.

In order that this may not occur, another formula must be explained.

There is:
\[
 \sqrt{x}=2\sqrt[4]{q}\,{1+q^2\over 1+2q}
 \qquad \text{fund. pag. 184}
\]
\[
 q={x^2\over16}\left\{{1+2q\over 1+q^2}\right\}^{4}
\]
\[
 \log q=\log{x^2\over16}-4\log(1+q^2)+4\log(1+2q)
 =\log{x^2\over16}+8q-12q^2+{8\over3}q^3.
\]
If $q$ is to be computed more accurately, it is necessary to compute some quantity $q_1$ approaching $q$ more closely, from the formula
\[
 \log q_1=\log{x^2\over16}+{x^2\over2};
\]
after this has been done, $q$ is defined from the equation
\[
 \log q=\log{x^2\over16}+8q_1-12q_1^2+{8\over3}q_1^3.
\]
\endgroup
